\documentclass[11pt]{article}

\usepackage[utf8]{inputenc}
\usepackage[T1]{fontenc}
\usepackage{lmodern}
\usepackage[margin=1in]{geometry}
\usepackage{amsmath, amssymb, amsthm, mathtools}
\usepackage{booktabs}
\usepackage{longtable}
\usepackage{array}
\usepackage{enumitem}
\usepackage[htt]{hyphenat}
\usepackage{hyperref}
\usepackage{xcolor}
\usepackage{graphicx}
\usepackage{float}

\hypersetup{
  colorlinks=true,
  linkcolor=blue!50!black,
  urlcolor=blue!50!black,
  citecolor=blue!50!black
}

\setlist[itemize]{topsep=4pt,itemsep=2pt}
\setlist[enumerate]{topsep=4pt,itemsep=2pt}
\setlength{\emergencystretch}{3em}
\newcolumntype{L}[1]{>{\raggedright\arraybackslash}p{#1}}

\title{From Listed Option Prices to an Arbitrage-Free Surface,\\Dupire Local Volatility, Pricing, and Hedging\\[4pt]
\large Theory and Repository Walkthrough}
\author{}
\date{\today}

\begin{document}
\sloppy

\maketitle
\tableofcontents
\newpage

\section{Purpose of This Document}

This document explains the full intellectual arc of the project in this repository. The intended reader already knows what an option is and has heard of the Black-Scholes equation, but does not yet have a working mental model of:

\begin{itemize}
  \item implied volatility,
  \item volatility surfaces,
  \item static arbitrage constraints across strike and maturity,
  \item Dupire local volatility,
  \item PDE pricing under local volatility,
  \item or how a real codebase turns raw listed option quotes into a hedging study.
\end{itemize}

The document therefore has two goals:

\begin{enumerate}
  \item explain the financial theory in a pedagogical order;
  \item explain exactly how this repository implements that theory.
\end{enumerate}

After reading it, you should be able to answer the main questions someone might ask about the project:

\begin{itemize}
  \item What problem is the project solving?
  \item Why is a volatility surface needed instead of one volatility number?
  \item What does ``arbitrage-free'' mean here?
  \item Why build a call-price surface and not just fit implied vol directly?
  \item What is Dupire local volatility, and why is it numerically delicate?
  \item How are options priced in the repository?
  \item How is the Black-Scholes versus Local Vol hedge study designed?
  \item What are the strengths and limitations of the empirical results?
\end{itemize}

\section{The Big Picture}

The project takes daily SPY option-chain snapshots and turns them into a sequence of increasingly structured objects:

\begin{enumerate}
  \item raw listed option quotes,
  \item a cleaned and filtered quote set,
  \item a smooth total-variance surface,
  \item an arbitrage-repaired call-price grid,
  \item an implied-volatility surface backed by that call grid,
  \item a Dupire local-volatility surface,
  \item numerical prices and deltas under Black-Scholes and Local Vol,
  \item and finally a dated historical hedging comparison on fixed listed contracts.
\end{enumerate}

The most important design choice is this:

\begin{quote}
The canonical object of record is not the raw fitted volatility function. It is the repaired call-price grid after projection onto the static no-arbitrage set.
\end{quote}

That choice matters because Dupire differentiation is very sensitive to noise. If the input call surface is inconsistent or economically impossible, the local-volatility surface is not trustworthy.

\section{Financial Background}

\subsection{From Option Prices to Risk-Neutral Information}

A European call option with strike $K$ and maturity $T$ pays
\[
(S_T - K)^+ = \max(S_T - K, 0)
\]
at expiration. A put pays
\[
(K - S_T)^+.
\]

If one knew the full risk-neutral distribution of $S_T$, then call prices for all strikes at that maturity could be computed by expectation under discounting. In practice we observe the reverse problem: the market gives us call and put prices, and we want to infer a consistent object that explains them.

This is why the project is really a surface construction problem. A single maturity gives one smile or skew. A collection of maturities gives a two-dimensional surface.

\subsection{Why Black-Scholes Is Not Enough}

In Black-Scholes, the underlying follows a geometric Brownian motion with constant volatility $\sigma$:
\[
dS_t = (r-q)S_t\,dt + \sigma S_t\,dW_t,
\]
where:

\begin{itemize}
  \item $r$ is the continuously compounded risk-free rate,
  \item $q$ is the continuous dividend yield,
  \item and $\sigma$ is constant.
\end{itemize}

That model leads to a closed-form price for European calls and puts. But in real markets, if you invert Black-Scholes option by option, the implied volatility depends on strike and maturity. That is the empirical smile or skew.

So Black-Scholes remains useful, but not as a literal description of the market. In modern practice it is often used as a quoting convention: the market price is converted into an implied volatility.

\subsection{Implied Volatility}

For a given option price $C^{\text{mkt}}$, implied volatility is the value $\sigma_{\text{imp}}$ such that
\[
C^{\text{BS}}(S_0,K,T,r,q,\sigma_{\text{imp}})=C^{\text{mkt}}.
\]

This repository computes that inversion with:

\begin{itemize}
  \item \texttt{implied\_vol()} and \texttt{implied\_vol\_one()} in \texttt{src/implied\_volatility/iv\_solver.py},
  \item using Black-Scholes pricing and vega from \texttt{src/implied\_volatility/black\_scholes.py}.
\end{itemize}

The solver supports a hybrid Newton-Brent approach because some options are easy to invert and others are not. Deep in/out-of-the-money options, very short maturities, and quotes near no-arbitrage bounds can all be numerically awkward.

\subsection{Why Total Variance Is Often Better Than Raw IV}

The repository does not treat implied vol itself as the smooth object to interpolate. It instead works with total variance
\[
w(x,T)=\sigma_{\text{imp}}(x,T)^2\,T.
\]

Why?

\begin{itemize}
  \item Total variance often behaves more smoothly across maturity than raw IV.
  \item Many surface parameterizations and interpolation schemes are easier to stabilize in total variance.
  \item Dupire is sensitive to curvature, so using a smoother state variable is helpful.
\end{itemize}

The conversion functions are:

\begin{itemize}
  \item \texttt{total\_variance\_from\_iv()} in \texttt{src/surface\_fitting/spline.py},
  \item \texttt{iv\_from\_total\_variance()} in the same file.
\end{itemize}

\subsection{Why Use Log-Forward-Moneyness}

The project works in
\[
x=\log\left(\frac{K}{F(T)}\right),
\qquad
F(T)=S_0 e^{(r-q)T}.
\]

This coordinate is useful because it centers the smile relative to the forward rather than the spot alone. It is a more natural cross-maturity coordinate than raw strike when rates and carry matter.

The repository constructs these derived columns in
\path{add_derived_columns()} in \path{src/market_data/transforms.py}.

\section{Static Arbitrage and Why It Matters}

\subsection{What ``Arbitrage-Free'' Means Here}

The repository focuses on static no-arbitrage constraints across listed vanilla option prices. For calls, the important ones are:

\paragraph{Bounds}
\[
\max(S_0 e^{-qT}-Ke^{-rT},0)\le C(K,T)\le S_0 e^{-qT}.
\]

\paragraph{Monotonicity in strike}
As strike increases, call price should not increase:
\[
\frac{\partial C}{\partial K}\le 0.
\]

\paragraph{Convexity in strike}
Call prices should be convex in strike:
\[
\frac{\partial^2 C}{\partial K^2}\ge 0.
\]
This is closely tied to the risk-neutral density. In continuous theory, the Breeden-Litzenberger identity links the density to the second strike derivative of call price.

\paragraph{Monotonicity in maturity}
At the same strike, longer maturity should not reduce call value:
\[
\frac{\partial C}{\partial T}\ge 0,
\]
under the standard calendar interpretation used by the project.

\subsection{How the Repository Checks These Constraints}

These checks live in \texttt{src/arbitrage/static\_checks.py}:

\begin{itemize}
  \item \texttt{check\_static\_arbitrage\_calls()}
  \item \texttt{finite\_diff\_dK()}
  \item \texttt{finite\_diff\_d2K()}
  \item \texttt{finite\_diff\_dT()}
\end{itemize}

The central idea is simple: once a call-price grid is available, finite differences can diagnose whether the grid violates the economic shape restrictions above.

\subsection{Projection Onto the No-Arbitrage Set}

This repository does not merely check arbitrage. It explicitly repairs the call-price surface.

The core projection machinery is in \texttt{src/arbitrage/projection.py}:

\begin{itemize}
  \item \texttt{project\_call\_surface\_static\_arbitrage()}
  \item \texttt{\_repair\_row\_convexity\_and\_monotonicity()}
  \item \texttt{\_column\_pava\_nondecreasing()}
  \item \texttt{summarize\_calendar\_violations()}
\end{itemize}

Conceptually, the projection does this:

\begin{enumerate}
  \item start from a raw call-price grid produced by the fitted surface;
  \item repair each maturity slice so it is decreasing and convex in strike;
  \item repair the term structure so it is nondecreasing in maturity;
  \item iterate until the final grid passes the static checks.
\end{enumerate}

That final projected grid is the economic backbone of the repository.

\section{From Raw Listed Quotes to a Synthetic Call Surface}

\subsection{Why Raw Option Chains Are Messy}

Real option chains contain:

\begin{itemize}
  \item wide spreads,
  \item stale quotes,
  \item thinly traded strikes,
  \item maturities with too few usable points,
  \item and inconsistencies between calls and puts.
\end{itemize}

If one feeds this directly into Dupire, the derivatives with respect to strike and maturity become numerically unstable.

\subsection{Data Ingestion}

The live and historical option-chain source in the current repository is Theta. The key functions are in \texttt{src/market\_data/theta\_loader.py}:

\begin{itemize}
  \item \texttt{ThetaSnapshotConfig}
  \item \texttt{ensure\_theta\_terminal\_running()}
  \item \texttt{fetch\_theta\_option\_chain\_snapshot()}
\end{itemize}

The loader:

\begin{enumerate}
  \item queries Theta for quotes, open interest, and volume,
  \item merges them into one option-chain table,
  \item gets the underlying spot near the target snapshot time,
  \item attaches carry estimates,
  \item and returns a standardized frame of listed contracts.
\end{enumerate}

\subsection{Validation and Derived Columns}

Two modules are used very early in the pipeline:

\begin{itemize}
  \item \texttt{validate\_and\_clean()} in \texttt{src/market\_data/validators.py}
  \item \texttt{add\_derived\_columns()} in \texttt{src/market\_data/transforms.py}
\end{itemize}

Validation removes obviously unusable rows. Derived-column construction adds:

\begin{itemize}
  \item time to expiry,
  \item forward price,
  \item moneyness,
  \item log-moneyness,
  \item and related quantities used later in fitting and hedging.
\end{itemize}

\subsection{Why Merge OTM Puts and OTM Calls}

One of the most important practical design choices in the repository is that it does not build its surface from calls alone. Instead it forms a synthetic call data set:

\begin{itemize}
  \item on the left wing, use out-of-the-money puts,
  \item convert them into synthetic calls with put-call parity,
  \item on the right wing, use out-of-the-money calls directly.
\end{itemize}

The conversion uses
\path{put_call_parity_call_from_put()} in \path{src/implied_volatility/black_scholes.py}.

The actual quote preparation happens in
\path{_prepare_synthetic_call_quotes()} inside
\path{src/pipeline/calibration.py}.

Why this matters:

\begin{itemize}
  \item left-wing information is often better represented by puts,
  \item right-wing information is often better represented by calls,
  \item and a merged OTM view gives broader and more meaningful support for the eventual density and Dupire steps.
\end{itemize}

\section{The Surface-Building Pipeline in This Repository}

\subsection{The Central Calibration Entry Point}

The most important function in the codebase is
\path{calibrate_option_snapshot()} in
\path{src/pipeline/calibration.py}.

It takes one cleaned option snapshot and returns a
\texttt{SnapshotCalibrationArtifacts} object containing:

\begin{itemize}
  \item cleaned quotes,
  \item the fitted total-variance surface,
  \item the raw and repaired call-price grids,
  \item the arbitrage-free implied-volatility surface,
  \item the local-volatility surface,
  \item density diagnostics,
  \item and summary statistics.
\end{itemize}

The configuration object is \texttt{SnapshotCalibrationConfig}.

\subsection{Quote Weighting}

Inside the calibration pipeline, the function
\texttt{\_compute\_quote\_weights()} gives more influence to better quotes. It favors:

\begin{itemize}
  \item smaller relative spreads,
  \item higher volume,
  \item higher open interest,
  \item and OTM-source quality.
\end{itemize}

This is a practical compromise: the repository is not pretending every midpoint quote is equally trustworthy.

\subsection{Direct Total-Variance Interpolation}

The repository no longer uses SVI as the active canonical surface-building method. Instead it uses direct interpolation of total variance.

The relevant code is in \texttt{src/surface\_fitting/spline.py}:

\begin{itemize}
  \item \texttt{build\_total\_variance\_grid\_from\_slices()}
  \item \texttt{fit\_spline\_surface()}
  \item the \texttt{SplineSurface} class.
\end{itemize}

The logic is:

\begin{enumerate}
  \item for each maturity slice, smooth or interpolate total variance as a function of log-forward-moneyness;
  \item place all slices onto a common $x$-grid;
  \item interpolate across maturities onto a common $T$-grid;
  \item expose the resulting object as \texttt{SplineSurface}.
\end{enumerate}

This surface is not yet arbitrage-free in itself. It is a smooth interpolant used to generate a raw call-price grid.

\subsection{From Surface to Raw Call Grid}

The function
\texttt{call\_price\_grid\_from\_iv\_surface()} in
\texttt{src/implied\_volatility/surface.py}
converts the interpolated total-variance or IV surface into a call-price grid using Black-Scholes pricing.

This is conceptually important:

\begin{itemize}
  \item interpolation occurs in volatility-like space,
  \item but the economic constraints are imposed in price space.
\end{itemize}

\subsection{The Arbitrage-Free IV Surface}

Once the call grid has been repaired, the repository wraps it in
\texttt{ArbitrageFreeIVSurface} from
\texttt{src/implied\_volatility/surface.py}.

That class:

\begin{itemize}
  \item stores the arbitrage-free call grid,
  \item interpolates calls inside the grid,
  \item and inverts Black-Scholes to return implied vol when needed.
\end{itemize}

So the repository's notion of ``arbitrage-free IV surface'' really means:

\begin{quote}
an implied-volatility view backed by an arbitrage-free call-price surface.
\end{quote}

\section{Risk-Neutral Density}

\subsection{Breeden-Litzenberger Idea}

Under sufficient smoothness,
\[
\frac{\partial^2 C}{\partial K^2}
\]
is proportional to the risk-neutral terminal density. Intuitively:

\begin{itemize}
  \item monotonicity tells us calls get cheaper as strike increases,
  \item convexity tells us the slope changes in an economically meaningful way,
  \item and the second derivative extracts information about how probability mass is distributed across terminal outcomes.
\end{itemize}

In the repository, density diagnostics live in \texttt{src/arbitrage/density.py}:

\begin{itemize}
  \item \texttt{breeden\_litzenberger\_density()}
  \item \texttt{density\_normalization\_diagnostic()}
\end{itemize}

The density mass is not a philosophical detail. It is a practical way to test whether the strike grid is too truncated or the surface still looks unstable.

\section{Dupire Local Volatility}

\subsection{What Local Volatility Is}

Local volatility replaces constant $\sigma$ with a deterministic function of state and time:
\[
dS_t = (r-q)S_t\,dt + \sigma_{\text{loc}}(S_t,t) S_t\,dW_t.
\]

The appeal of Dupire's formula is that, in principle, if you know a smooth enough call surface $C(K,T)$ for all strikes and maturities, then you can infer a local variance surface:
\[
\sigma_{\text{loc}}^2(K,T)
=
\frac{
\frac{\partial C}{\partial T}
+(r-q)K\frac{\partial C}{\partial K}
+qC
}{
\tfrac12 K^2 \frac{\partial^2 C}{\partial K^2}
}.
\]

\subsection{Why Dupire Is Numerically Delicate}

This formula is attractive and dangerous at the same time:

\begin{itemize}
  \item it differentiates in maturity,
  \item it differentiates twice in strike,
  \item and the denominator can become tiny in wings or sparse regions.
\end{itemize}

That means any noise or residual inconsistency in the call surface gets amplified.

\subsection{How the Repository Extracts Dupire Vol}

The formula is implemented in \texttt{src/local\_volatility/dupire.py}:

\begin{itemize}
  \item \texttt{DupireConfig}
  \item \texttt{dupire\_local\_var\_from\_call\_grid()}
  \item \texttt{dupire\_local\_vol\_from\_call\_grid()}
\end{itemize}

The code uses finite differences and includes numerical floors and optional caps to avoid nonphysical blowups.

\subsection{Regularization}

Raw Dupire output is rarely good enough to use directly. This repository regularizes it in
\path{src/local_volatility/regularization.py}:

\begin{itemize}
  \item \texttt{LocalVolRegularizationConfig}
  \item \texttt{clip\_local\_vol()}
  \item \texttt{repair\_short\_end\_wings()}
  \item \texttt{smooth\_local\_vol\_gaussian()}
  \item \texttt{regularize\_local\_vol()}
\end{itemize}

The idea is not to hide problems. The idea is to make the local-vol surface usable while being explicit about where it is and is not trustworthy.

\subsection{Trusted and Supported Domains}

One mature feature of the repository is the domain mask. The local-vol surface is not treated as equally reliable everywhere.

The query object \texttt{LocalVolSurface} in
\texttt{src/local\_volatility/surface.py}
carries a discrete domain classification:

\begin{itemize}
  \item \texttt{0}: extrapolation or unsupported region,
  \item \texttt{1}: weakly supported region,
  \item \texttt{2}: trusted interior region.
\end{itemize}

This is operationally useful because it lets the project say:

\begin{quote}
the pipeline produced a local-vol surface, but only part of that surface is good enough for strong claims.
\end{quote}

That is much more honest than pretending the wings are as reliable as the liquid core.

\section{Pricing Under the Model}

\subsection{Black-Scholes Pricing}

Black-Scholes pricing and bounds live in
\texttt{src/implied\_volatility/black\_scholes.py}.

Important functions are:

\begin{itemize}
  \item \texttt{bs\_price()}
  \item \texttt{bs\_vega()}
  \item \texttt{bs\_delta()}
  \item \texttt{no\_arbitrage\_bounds\_call()}
  \item \texttt{no\_arbitrage\_bounds\_put()}
\end{itemize}

Black-Scholes plays three roles in the repository:

\begin{enumerate}
  \item it provides the implied-vol quoting convention,
  \item it provides a baseline pricing and hedging model,
  \item it converts the fitted volatility surface into raw call prices before projection.
\end{enumerate}

\subsection{Why a PDE Solver Is Needed}

Under local volatility, closed-form prices are usually not available. This repository therefore uses a Crank-Nicolson finite-difference PDE solver in
\texttt{src/pricing/pde\_solver.py}.

Key objects and functions:

\begin{itemize}
  \item \texttt{PDEConfig}
  \item \texttt{solve\_european\_pde\_local\_vol\_surface()}
  \item \texttt{price\_european\_pde\_local\_vol()}
\end{itemize}

Conceptually, the PDE solver:

\begin{itemize}
  \item lays down a spot grid,
  \item works backward from the terminal payoff,
  \item uses local volatility at each point on the grid,
  \item and returns option values $V(t,S)$ and eventually deltas.
\end{itemize}

\subsection{Monte Carlo}

The repository also includes local-vol Monte Carlo pricing in
\texttt{src/pricing/monte\_carlo.py}, primarily as an additional pricing engine and cross-check rather than as the main hedging engine.

That is useful because a serious derivatives project should not rely on one numerical method only.

\section{Delta Hedging}

\subsection{What Delta Hedging Means}

Delta hedging means dynamically trading the underlying so that the portfolio's first-order sensitivity to spot moves is offset.

If an option has delta $\Delta_t$, then a simple hedge holds $\Delta_t$ shares of the underlying against the short option position. In discrete time, the hedge is rebalanced whenever a new market observation arrives.

\subsection{Low-Level Hedging Engine}

The lower-level pathwise hedging engine is in \texttt{src/hedging/delta\_hedger.py}.

Key pieces:

\begin{itemize}
  \item \texttt{BSDeltaModel}
  \item \texttt{LocalVolPDEDeltaModel}
  \item \texttt{LocalVolPDEDeltaModelStrikeSurface}
  \item \texttt{DeltaHedger}
  \item \texttt{evolve\_cash()}
  \item \texttt{option\_payoff()}
\end{itemize}

The \texttt{DeltaHedger} class simulates:

\begin{enumerate}
  \item selling or buying the option at inception,
  \item financing the hedge in a cash account accruing at rate $r$,
  \item rebalancing the stock position at discrete times,
  \item charging transaction costs,
  \item and comparing the terminal portfolio to the realized option payoff.
\end{enumerate}

\subsection{Transaction Costs}

Transaction cost logic is in \path{src/hedging/transaction_costs.py}. The class \path{TransactionCostModel} supports:

\begin{itemize}
  \item no costs,
  \item proportional costs,
  \item fixed costs,
  \item and bid-ask style approximations.
\end{itemize}

The current canonical study uses explicit transaction costs instead of pretending hedging is free.

\section{The Historical Panel Backtest}

\subsection{Why the Repository Uses a Dated Panel}

One of the most important architectural changes in the project was to move away from the idea of freezing one surface and replaying it through unrelated historical spot paths. The current canonical design is a true dated-panel study:

\begin{itemize}
  \item each trade date has its own option snapshot,
  \item each trade date is recalibrated separately,
  \item the same fixed listed contract is then tracked over time,
  \item and Black-Scholes and Local Vol are compared on the same market liability.
\end{itemize}

\subsection{Contract Selection}

The dated-panel logic is implemented in
\texttt{src/hedging/market\_panel\_backtest.py}.

Important functions:

\begin{itemize}
  \item \texttt{select\_listed\_contracts()}
  \item \texttt{collect\_contract\_marks()}
  \item \texttt{build\_daily\_calibration\_cache()}
  \item \texttt{run\_daily\_market\_panel\_backtest()}
\end{itemize}

The contract selector deliberately focuses on a liquid short-dated core:

\begin{itemize}
  \item near-ATM contracts,
  \item tighter spread filters,
  \item minimum volume and open interest thresholds,
  \item and short target maturities.
\end{itemize}

That narrower scope is a feature, not a defect. With retail-access data, fitting the liquid core well is more credible than pretending the whole chain is equally high quality.

\subsection{How Quotes Are Generated for Each Model}

Inside \texttt{market\_panel\_backtest.py}, two helper functions matter:

\begin{itemize}
  \item \texttt{\_bs\_quote()}
  \item \texttt{\_lv\_quote()}
\end{itemize}

\texttt{\_bs\_quote()}:

\begin{itemize}
  \item takes the arbitrage-free IV surface,
  \item evaluates the implied volatility at the contract's strike and remaining maturity,
  \item then produces a Black-Scholes price and delta.
\end{itemize}

\texttt{\_lv\_quote()}:

\begin{itemize}
  \item takes the calibrated local-vol surface for that trade date,
  \item solves a PDE on a spot grid around the contract,
  \item and extracts the option price and delta from the PDE surface.
\end{itemize}

\subsection{How Results Are Summarized}

Performance metrics are summarized in
\texttt{src/hedging/metrics.py}, especially:

\begin{itemize}
  \item \texttt{summarize\_backtest()}
  \item \texttt{summarize\_market\_panel\_backtest()}
\end{itemize}

The project reports multiple metrics because no single error number tells the whole story:

\begin{itemize}
  \item entry pricing error,
  \item replication error,
  \item net market PnL,
  \item transaction costs,
  \item number of rebalances,
  \item mean absolute mark-to-market error.
\end{itemize}

This is deliberate. A model can look bad because it priced badly at inception, because it hedged badly after inception, or both. The repository tries not to collapse those effects into one opaque number.

\section{Operational Scripts and How the Repository Runs Day to Day}

\subsection{Daily Snapshot Collection}

The script
\texttt{scripts/archive\_option\_snapshot.py}
collects one dated option snapshot and archives it. It is the lowest-level operational entry point for daily collection.

\subsection{Backfill}

The script
\texttt{scripts/backfill\_theta\_option\_panel.py}
uses historical Theta access to backfill a month or more of 15:45 snapshots into the canonical panel.

\subsection{Single-Snapshot Calibration}

The script
\texttt{scripts/run\_snapshot\_calibration.py}
runs calibration for one specific date or one specific snapshot file and writes the resulting artifacts.

\subsection{Daily Canonical Workflow}

The main orchestration script is
\texttt{scripts/run\_canonical\_daily\_update.py}.

Its job is to:

\begin{enumerate}
  \item load the daily configuration,
  \item fetch or reuse the latest snapshot,
  \item append the panel,
  \item clear stale current-output artifacts,
  \item calibrate the latest date,
  \item run the market-panel backtest,
  \item generate the report and figures,
  \item and archive the current run as a milestone snapshot.
\end{enumerate}

\subsection{Backtest and Reporting Scripts}

Two more scripts matter:

\begin{itemize}
  \item \texttt{scripts/run\_backtest\_market\_panel.py}
  \item \texttt{scripts/generate\_report.py}
\end{itemize}

The first runs the dated-panel hedge study. The second renders the figures and summary markdown from the generated artifacts.

\section{How to Explain the Repository in a Conversation}

\subsection{One-Sentence Summary}

\begin{quote}
This repository builds a daily pipeline from listed SPY option chains to an arbitrage-free implied-volatility surface, derives a Dupire local-volatility model from that surface, prices with PDE and Monte Carlo, and compares Black-Scholes versus Local Vol hedging on fixed listed contracts with transaction costs.
\end{quote}

\subsection{A Slightly Longer Summary}

\begin{quote}
The project starts from real listed option quotes, cleans them, merges OTM puts and calls into a synthetic call surface, fits total variance in log-forward-moneyness and maturity, projects the resulting call grid onto the static no-arbitrage set, extracts Dupire local vol, prices numerically, and then runs a dated-panel historical hedge study in which both Black-Scholes and Local Vol hedge the same listed contracts using the same market entry premiums.
\end{quote}

\section{What Is Strong About the Project}

\begin{itemize}
  \item It is end to end.
  \item It does not stop at implied-vol fitting; it goes through Dupire, pricing, and hedging.
  \item It treats no-arbitrage as an explicit numerical constraint, not just a plotting nicety.
  \item It uses real listed data and a dated panel rather than only synthetic examples.
  \item It explicitly models transaction costs.
  \item It distinguishes trusted and untrusted parts of the local-vol surface.
  \item It has operational scripts, tests, and documentation rather than living only in notebooks.
\end{itemize}

\section{What Its Current Limitations Are}

\begin{itemize}
  \item The stock-side inputs still use a small Yahoo fallback.
  \item Even with better option data, calibration fit to the market is not perfect.
  \item The trusted local-vol region is narrower than the full numerical grid.
  \item Local Vol does not simply dominate Black-Scholes in the live empirical results.
  \item SPY may not be the asset on which deterministic local volatility is most naturally superior.
\end{itemize}

These are not fatal flaws. They are the normal limits of a serious research project built on noninstitutional infrastructure.

\section{Common Questions and Good Answers}

\subsection*{Why not just use Black-Scholes?}

Black-Scholes is still used, but the market does not quote one volatility for all strikes and maturities. The whole point of the surface is to explain the observed smile and term structure.

\subsection*{Why not fit implied volatility directly and stop there?}

Because the eventual goal is pricing and hedging under a dynamic model. The IV surface is a static description of today's prices. Dupire turns that into a state-dependent diffusion model.

\subsection*{Why work so much in price space?}

Because no-arbitrage constraints are economically natural in call prices. A volatility surface can look smooth while still implying inconsistent prices. Projecting in price space is therefore the safer anchor.

\subsection*{Why merge OTM puts and calls?}

Because each wing is usually better informed by the instrument that is actually out of the money there. This gives a more meaningful synthetic call surface than relying on calls alone.

\subsection*{Why is local volatility hard in practice?}

Because Dupire differentiates a noisy surface, especially with respect to strike. That amplifies data noise and interpolation error, so cleaning, projection, and regularization all matter.

\subsection*{Why a PDE solver?}

Local Vol generally does not admit a closed-form option price, so a finite-difference PDE solver is one standard numerical method to compute prices and deltas.

\subsection*{Why does the repository not claim Local Vol wins everywhere?}

Because the project is trying to be honest. The empirical result is mixed, and that is informative. A strong project is not one that declares victory regardless of the data.

\section{Module-by-Module Walkthrough}

\renewcommand{\arraystretch}{1.2}
\begin{longtable}{L{0.28\textwidth} L{0.65\textwidth}}
\toprule
\textbf{File} & \textbf{Role in the pipeline} \\
\midrule
\endfirsthead
\toprule
\textbf{File} & \textbf{Role in the pipeline} \\
\midrule
\endhead
\path{src/market_data/theta_loader.py} & Fetches daily or historical option-chain snapshots from Theta and attaches underlying spot and carry inputs. \\
\path{src/market_data/validators.py} & Removes clearly unusable option rows before surface fitting. \\
\path{src/market_data/transforms.py} & Adds forward, time-to-expiry, moneyness, and log-moneyness coordinates. \\
\path{src/pipeline/calibration.py} & Core single-snapshot calibration entry point. Builds the synthetic call data, fits total variance, projects the call grid, computes density, extracts local vol, and packages artifacts. \\
\path{src/surface_fitting/spline.py} & Canonical direct total-variance interpolation machinery. \\
\path{src/arbitrage/static_checks.py} & Finite-difference diagnostics for bounds, strike monotonicity, strike convexity, and calendar monotonicity. \\
\path{src/arbitrage/projection.py} & Repairs the raw call grid into an arbitrage-free grid. \\
\path{src/arbitrage/density.py} & Computes density diagnostics from the repaired call surface. \\
\path{src/implied_volatility/black_scholes.py} & Black-Scholes pricing, delta, vega, parity relations, and no-arbitrage bounds. \\
\path{src/implied_volatility/iv_solver.py} & Inverts Black-Scholes to compute implied volatilities. \\
\path{src/implied_volatility/surface.py} & Wraps the arbitrage-free call grid as an implied-volatility surface object. \\
\path{src/local_volatility/dupire.py} & Applies Dupire's formula to the repaired call grid. \\
\path{src/local_volatility/regularization.py} & Stabilizes the raw local-vol surface. \\
\path{src/local_volatility/surface.py} & Query object for local volatility, including trusted-domain classification. \\
\path{src/pricing/pde_solver.py} & Crank-Nicolson PDE solver for local-vol pricing. \\
\path{src/pricing/monte_carlo.py} & Monte Carlo pricing under local volatility. \\
\path{src/hedging/delta_hedger.py} & Low-level delta-hedging mechanics for one realized path. \\
\path{src/hedging/market_panel_backtest.py} & High-level dated-panel backtest using fixed listed contracts and daily recalibration. \\
\path{src/hedging/transaction_costs.py} & Transaction-cost model used by the hedge engine. \\
\path{src/hedging/metrics.py} & Summary statistics for backtests. \\
\path{scripts/archive_option_snapshot.py} & Archives one dated snapshot. \\
\path{scripts/backfill_theta_option_panel.py} & Backfills historical 15:45 snapshots into the panel. \\
\path{scripts/run_snapshot_calibration.py} & Runs one calibration job and writes artifacts. \\
\path{scripts/run_backtest_market_panel.py} & Runs the dated-panel historical hedge study. \\
\path{scripts/run_canonical_daily_update.py} & End-to-end daily automation. \\
\path{scripts/generate_report.py} & Builds the current report and figures from artifacts. \\
\bottomrule
\end{longtable}

\section{Final Mental Model}

The right way to think about this repository is:

\begin{enumerate}
  \item \textbf{Data problem:} listed option chains are noisy and inconsistent.
  \item \textbf{Surface problem:} one needs a smooth, arbitrage-free representation across strike and maturity.
  \item \textbf{Model problem:} Dupire local volatility translates that surface into a state-dependent diffusion.
  \item \textbf{Numerical problem:} pricing and delta extraction require PDE and related numerical machinery.
  \item \textbf{Evaluation problem:} a model should be judged not only by fit, but also by hedging behavior on dated listed contracts with explicit costs.
\end{enumerate}

This repository addresses all five.

\end{document}
